\documentclass[11pt,reqno]{amsart}
\usepackage{amsmath,amssymb,mathtools}
\usepackage[italian]{babel}
\usepackage{microtype}
\usepackage{tikz}
\usepackage{placeins}
\usetikzlibrary{arrows.meta,positioning,calc,decorations.pathreplacing}
\usepackage[letterpaper,textwidth=6.25in,textheight=8.85in,centering]{geometry}
\usepackage[hidelinks,unicode]{hyperref}
\usepackage{booktabs,array,longtable}
\usepackage{pgfplots}
\pgfplotsset{compat=1.18}
\hypersetup{pdftitle={Pulizia chimica del grafene trasferito con PMMA e PPC},pdfauthor={Codex, su richiesta di Mattia Apicella},pdfsubject={Revisione 6: pressione corretta, campionamento molecolare e autoassociazione dei solventi}}
\numberwithin{equation}{section}
\renewcommand{\abstractname}{Sommario}
\renewcommand{\refname}{Riferimenti}
\renewcommand{\keywordsname}{Parole chiave}
\renewcommand{\datename}{\textit{Data}:}
\newcommand{\SiO}{\ensuremath{\mathrm{SiO}_2}}
\newcommand{\HCl}{\ensuremath{\mathrm{HCl}}}
\newcommand{\degC}{\ensuremath{^{\circ}\mathrm{C}}}
\newcommand{\um}{\ensuremath{\mu\mathrm{m}}}
\newcommand{\nm}{\ensuremath{\mathrm{nm}}}

\title[Grafene: teoria e calcoli molecolari della pulizia]{Pulizia chimica del grafene gi\`a trasferito con PMMA/PPC:\\ teoria, calcoli molecolari e verifica dei modelli}
\author{Codex}
\date{6 ottobre 2026 --- revisione 6}
\thanks{Studio redatto su richiesta di Mattia Apicella per il confronto con Armando Serio. Include calcoli eseguiti e verificati, ma nessun nuovo esperimento sui materiali. Le costanti cinetiche delle simulazioni sono scenari dichiarati, non stime del campione. Questa revisione sostituisce la raccomandazione iniziale di HCl alla luce dei tentativi gi\`a effettuati.}
\keywords{Grafene CVD, PMMA AR-P 672.045, PPC, desorbimento, sequenza dei solventi, bilanci di massa}
\begin{document}
\begin{abstract}
Il problema considerato \`e rimuovere chimicamente residui da grafene
su \SiO/Si, trasferito dal rame con PMMA AR-P 672.045 e PPC.
L'insuccesso riferito di numerosi solventi, acidi e prodotti commerciali
esclude di ripresentarli come nuove soluzioni. L'unico trattamento con
un modesto effetto combina DMF, isopropanolo e THF. Proponiamo di
verificare se il risultato dipenda dalla preparazione di uno stato
polimerico estraibile, dalla sua successiva perdita durante i
risciacqui, oppure da una frazione persistente di diversa natura.
Deriviamo un modello conservativo a tre stati che mostra quando
l'ordine dei liquidi conta e quando l'alternanza aiuta o peggiora.
Un secondo modello separa distacco, trasporto e riadsorbimento;
un bilancio indipendente limita la massa depositabile durante
l'asciugatura. Le simulazioni includono controesempi e controlli
analitici. Ne segue una proposta concreta: confrontare la sequenza
esistente con varianti prive dei passaggi IPA e, se emerge una
mobilizzazione utile, con alternanze a esposizioni complessive uguali.
L'obiettivo \`e amplificare un meccanismo misurabile conservando il
grafene. Aggiungiamo calcoli quantistici, due modelli polarizzabili appresi e traiettorie di un oligomero su grafene in cinque liquidi. Il confronto con un acido gi\`a fallito mostra che una forte interazione locale con l'estere non basta a scegliere il pulente. Il controllo del liquido impedisce di promuovere un modello dalla sola densit\`a. L'efficacia sul campione reale rimane da verificare.
\end{abstract}
\maketitle

\section{La proposta che sceglierei ora}

\textbf{La mia ipotesi di lavoro \`e che la frazione ancora rimovibile
sia limitata dal passaggio tra stati interfaciali del polimero e dalla
successiva estrazione, pi\`u che dalla scelta di un altro solvente
generico.} La prova pi\`u immediata conserva DMF e THF, ma separa
l'effetto dell'IPA intermedio da quello finale. La prova successiva
confronta quattro brevi alternanze con gli stessi liquidi raggruppati
in blocchi, mantenendo uguali esposizioni e ricambi.

Non assumo in partenza che l'IPA sia dannoso. Pu\`o favorire la
precipitazione o il riattacco di materiale gi\`a mobilizzato, oppure
modificare la conformazione di una frazione aderente e renderla
estraibile nel bagno seguente. Queste ipotesi danno previsioni
opposte. Il modello delle Sezioni~\ref{sec:states}--\ref{sec:sim}
le rende quantitative senza attribuire al campione costanti che
nessuno ha misurato.

La modifica da provare per prima, se occorre scegliere un solo
nuovo percorso, \`e \emph{DMF con ricambi freschi, seguito direttamente
da THF, mantenendo inizialmente invariato il risciacquo finale}.
In questo modo si verifica se l'IPA intermedio annulli una
mobilizzazione ottenuta nel DMF. La Tabella~\ref{tab:factorial}
specifica tempi e controlli. Se invece quel passaggio risulta utile,
si passa alla prova di alternanza della Tabella~\ref{tab:cycles}.

La scelta \`e motivata dall'unico percorso che ha gi\`a prodotto un
segnale favorevole. Non \`e deducibile una probabilit\`a di successo
dal suo semplice esistere: il miglioramento riferito non \`e ancora
quantificato e i confronti precedenti non sono descritti abbastanza
da isolare l'effetto della sequenza.


\paragraph{Aggiornamento molecolare.}
HFIP rimane una candidata da verificare, discussa nella Sezione~\ref{sec:molecular}. I nuovi controlli non la promuovono sopra DMF.
Esiste evidenza di dissoluzione del PMMA in HFIP a temperatura ambiente,
ma la sua efficacia sui residui aderenti del caso non \`e dimostrata.
I calcoli eseguiti hanno anzi scartato il criterio semplicistico
``interazione pi\`u forte con il carbonile, quindi pulizia migliore''.
La sequenza DMF--THF rimane il controllo operativo; una classifica
quantitativa fra DMF, THF e HFIP richiede ancora un modello interfaciale
qualificato, oltre alla verifica sul materiale.

\section{Il caso reale e i risultati negativi da conservare}

Il materiale iniziale indicato \`e PMMA AR-P 672.045, che il produttore
identifica come 950K al 4,5\% in anisolo \cite{Allresist}. Il campione
comprende anche PPC sopra il PMMA, senza contatto intenzionale del
PPC con il grafene. Sono riferiti un passaggio a 90\,\degC\ per
2 minuti, distacco elettrochimico dal Cu e trasferimento su \SiO/Si.
Elettrolita, spessori, eventuali altre cotture e identificazione
chimica delle particelle non sono ancora noti.
Le particelle sono riferite subito dopo il trasferimento, prima della
litografia. La membrana PC e il supporto PDMS sono impiegati in un
successivo prelievo a secco: questo uso futuro non identifica l'origine
dei depositi gi\`a presenti. La stima ``90\% PMMA'' comunicata dall'esperto
\`e una valutazione di probabilit\`a, non una misura della composizione.

Il lavoro di Tyagi, Mi\v{s}eikis, Coletti e collaboratori
\cite{Tyagi2022} corrisponde bene a questo schema PMMA/PPC e ai
passaggi termici indicati. Non attribuiamo automaticamente al caso
attuale tutti i parametri del loro supplemento. Quel lavoro \`e
utile come riferimento di processo; il suo solvente commerciale
non \`e una nuova risposta all'elenco di tentativi di Armando.

Sono gi\`a stati provati, secondo quanto comunicato:
\begin{itemize}
\item HCl per 48 ore e acido acetico glaciale;
\item acetato di etile, cloroformio, DMF, THF, MPK, acetone;
\item IPA/acqua 75/25 e prodotti indicati come AR600.7 e AR300-76,
quest'ultimo sia a temperatura ambiente sia a 80\,\degC.
\end{itemize}
Manteniamo ``MPK'' e ``AR600.7'' come sigle riferite, senza
correggerle arbitrariamente in altri prodotti. Il precedente con
HCl \cite{Xiao2019} resta un risultato pubblicato, ma \textbf{HCl
48 ore viene ritirato come nuova raccomandazione per questo caso}.
La concentrazione mancante limita un confronto quantitativo, senza
cancellare il risultato negativo ricevuto.

L'unica pulizia descritta come debolmente efficace \`e
\begin{equation}\label{eq:real}
\begin{split}
 &\mathrm{DMF}\ (30\,\mathrm{min})
 \longrightarrow \mathrm{IPA}\ (5+3\,\mathrm{min})\\[-2pt]
 &\longrightarrow \mathrm{THF}\ (30\,\mathrm{min})
 \longrightarrow \mathrm{IPA}\ (5+3\,\mathrm{min})
 \longrightarrow \mathrm{N_2}.
\end{split}
\end{equation}
La somma dei bagni \`e 76 minuti. L'agitazione riferita \`e 500 rpm:
non descriviamo quindi il trattamento come un bagno quiescente.
Questa velocit\`a non determina per\`o il moto locale al campione
senza geometria, posizione, volume e tipo di agitatore.

\section{Che cosa aggiunge la letteratura estesa}

\subsection{La sequenza pu\`o cambiare il comportamento del polimero}
Manjkow e collaboratori \cite{Manjkow1987} osservano mediante
ellissometria un passaggio stretto fra dissoluzione e rigonfiamento
del PMMA in miscele solvente/non solvente. Le miscele studiate non
sono quelle di Armando: il lavoro giustifica studiare la traiettoria
di composizione, non applicare qui le sue soglie. Rooks e
collaboratori \cite{Rooks2002} mostrano inoltre perch\'e IPA e
IPA/acqua non si possano classificare con una sola etichetta
indipendente da massa molecolare e stato del resist.

La solubilit\`a del polimero libero non determina il distacco
dall'interfaccia. Simulazioni di desorbimento di catene adsorbite
\cite{Huston2020} trovano regimi cinetici molto diversi al variare
dell'adsorbimento. Si tratta di modelli generici, non di una
parametrizzazione di AR-P 672.045 su grafene. Non \`e corretto
ricavare una barriera moltiplicando l'energia di un monomero per
tutta la catena: distacco progressivo, conformazione e contatti
effettivi contano.

\subsection{Perch\'e il semplice ricambio non \`e gi\`a la soluzione}
Ayodele e collaboratori \cite{Ayodele2022} impiegano PMMA 950K,
distacco elettrochimico dal rame e un estrattore Soxhlet con acetone
distillato. Osservano miglioramenti morfologici, ma il confronto
varia anche temperatura, durata e risciacquo; manca una misura XPS
che certifichi l'assenza del residuo. La storia termica \`e diversa
da quella comunicata per questo caso.

Il risultato del 2025 di Tang e collaboratori \cite{Tang2025}
\`e particolarmente istruttivo: un processo ciclico migliora
l'aspetto macroscopico, ma i rapporti XPS
C=C:C--O:O--C=O sono circa 100:28:8 nel controllo e 100:30:7
nel processo ciclico. I residui rimangono. Questo impedisce di
promettere che solvente rinnovato o assenza di IPA finale bastino
da soli. Il lavoro non testa l'alternanza DMF/IPA/THF qui proposta.

\subsection{Miscela IPA/acqua, etanolo e prestazioni elettriche}
Duan e collaboratori \cite{Duan2022} riportano un miglioramento
con IPA/acqua 75/25, ma dopo acetone e ricottura in vuoto a
200\,\degC. La loro frazione di componenti C--O e C=O resta
misurabile; inoltre il controllo in sola acqua subisce danni.
Il metodo \`e gi\`a nell'elenco di Armando e non viene riproposto.

Merino e collaboratori \cite{Merino2024} studiano residui da
fabbricazione anche con etanolo e THF: documentano miglioramenti
elettrici e chimici parziali, ma anche problemi di distacco con
THF. Non \`e una prova sul nostro PMMA/PPC. Suk e collaboratori
\cite{Suk2013} forniscono un controesempio ancora pi\`u netto:
formammide assorbita nei residui migliora le propriet\`a elettriche
senza corrispondente scomparsa della morfologia residua.
Una migliore curva del transistor, da sola, non decide la pulizia.

\subsection{Acido, identit\`a chimica e diagnosi del residuo}
Nel PMMA l'idrolisi dell'estere laterale non rompe automaticamente
la catena carboniosa. Nei polimetacrilati di Sch\"onemann e
collaboratori \cite{Schoenemann2018}, diversi dal PMMA, la stabilit\`a
all'acido mostra quanto sia rischioso inferire la cinetica dal solo
gruppo funzionale. Anche il PPC ha una risposta all'idrolisi
dipendente dalle condizioni \cite{Jung2006}. Questi studi non
forniscono una nuova miscela di pulizia validata per il campione.

Wang e collaboratori \cite{Wang2017} usano marcatura isotopica del
PMMA per seguirne i residui. Schwartz e collaboratori
\cite{Schwartz2019} identificano localmente polimeri intrappolati
in eterostrutture mediante spettroscopia infrarossa associata
all'AFM. Non studiano il nostro PPC: il loro PC \`e un materiale
diverso. Il punto trasferibile \`e il metodo di identificazione.
Una particella vista in AFM non acquista identit\`a chimica dalla
sola somiglianza con un residuo di PMMA.

\section{Una teoria della sequenza: preparare, estrarre, ribloccare}
\label{sec:states}

\subsection{Una previsione che non richiede conoscere i tassi}
Supponiamo che ogni liquido $j$ rimuova una stessa popolazione di
residuo $m$ con un tasso costante $k_j$, senza memoria:
\begin{equation}\label{eq:scalar}
 \dot m=-k_jm.
\end{equation}
Dopo tempi totali $T_j$ nei vari liquidi,
\begin{equation}\label{eq:orderindependent}
 m_{\mathrm{finale}}=m_0\exp\!\left(-\sum_j k_jT_j\right).
\end{equation}
L'ordine e il numero di alternanze non compaiono. Dunque una
differenza riproducibile fra ordine a blocchi e alternato, a pari
esposizioni, ricambi e condizioni, \emph{falsifica questo modello}.
Non identifica da sola il meccanismo alternativo: concentrazione
del bagno, trasformazioni, stati conformazionali e danno meccanico
possono introdurre memoria. L'osservazione finora riferita non \`e
ancora un confronto controllato di questo tipo.

\subsection{Il modello minimo che permette un effetto dell'ordine}
Dividiamo la massa in tre stati:
$L$, residuo ancora aderente e non prontamente estraibile;
$M$, residuo mobilizzato ma ancora sul campione;
$R$, materiale gi\`a esportato fuori dalla zona di riadesione.
Nel liquido $j$ poniamo
\begin{equation}\label{eq:generator}
 \frac{d}{dt}
 \begin{pmatrix}L\\M\\R\end{pmatrix}
 =K_j\begin{pmatrix}L\\M\\R\end{pmatrix},\qquad
 K_j=\begin{pmatrix}
 -a_j&b_j&0\\a_j&-(b_j+e_j)&0\\0&e_j&0
 \end{pmatrix}.
\end{equation}
Il tasso $a_j$ mobilizza, $b_j$ riblocca, $e_j$ esporta.
Sono non negativi e hanno dimensione tempo$^{-1}$.
La somma $L+M+R$ si conserva e masse inizialmente non negative
rimangono tali. L'assorbimento nello stato $R$ presuppone un
allontanamento efficace: se il polimero resta nel bagno e torna
sulla superficie, occorre aggiungere il comparto della
Sezione~\ref{sec:transport}.

Per due trattamenti D e T, con tempi $t_D,t_T$, l'ordine D poi T
d\`a
\begin{equation}\label{eq:ordered}
 v_{DT}=e^{K_Tt_T}e^{K_Dt_D}v_0,
 \qquad v=(L,M,R)^{\mathsf T}.
\end{equation}
I due esponenziali generalmente non commutano. Per tempi piccoli,
la differenza di ordine parte da
\begin{equation}\label{eq:commutator}
 v_{DT}-v_{TD}=(K_TK_D-K_DK_T)v_0\,t_Dt_T+
 O\bigl((t_D+t_T)^3\bigr).
\end{equation}
Questa \`e una ragione quantitativa per cui due solventi
mediocri separatamente potrebbero essere complementari. Non
assegniamo automaticamente a DMF il ruolo $a$ e a THF il ruolo
$e$: sono proprio i ruoli da mettere alla prova.

\subsection{Lo stesso IPA pu\`o aiutare o ostacolare}
Consideriamo un caso esatto: nel passaggio P non vi \`e esportazione,
ma $L\rightleftarrows M$ con tassi $a,b$; nel successivo G vi sono
$M\to R$ a tasso $k$ e $M\to L$ a tasso $c$. Per tempi $t,u$,
la massa esportata in pi\`u da P poi G rispetto a G poi P \`e
\begin{equation}\label{eq:sign}
 \Delta R=
 \frac{k}{k+c}\bigl(1-e^{-(k+c)u}\bigr)
 \bigl(1-e^{-(a+b)t}\bigr)
 \frac{aL_0-bM_0}{a+b}.
\end{equation}
La formula vale per $a+b>0$ e $k+c>0$, con i casi nulli ottenuti
per continuit\`a. Quando i prefattori sono positivi, il segno
\`e quello di $aL_0-bM_0$.

Se la superficie \`e soprattutto bloccata e P la mobilizza,
il passaggio intermedio aiuta. Se il DMF ha gi\`a prodotto molta
frazione $M$ e l'IPA la riblocca, lo stesso passaggio ostacola
l'estrazione. Questa previsione \`e pi\`u precisa di ``l'IPA fa
precipitare'' o ``l'IPA pulisce''. L'ordine inverso termina in
un liquido diverso: la verifica deve misurare materiale esportato
prima dell'asciugatura o imporre un passaggio finale comune.

\subsection{Perch\'e la contrazione non \`e automaticamente un vantaggio}
Un solvente meno favorevole pu\`o ridurre il volume occupato dalla
catena, ma anche aumentare aggregazione e adesione alla superficie.
Per rendere utile un ciclo occorrono due condizioni: una quota
misurabile passi a uno stato pi\`u estraibile, e il liquido seguente
la allontani prima che torni aderente. Una diminuzione dell'impronta
AFM con aumento dell'altezza e volume conservato suggerisce
riorganizzazione, non ancora rimozione. Il bagno seguente deve
ridurre la massa o il volume con un riscontro chimico.

\section{Calcoli eseguiti e loro conseguenze}
\label{sec:sim}

Abbiamo calcolato gli esponenziali dei generatori con una somma
di matrici stocastiche, detta uniformizzazione, evitando tassi
negativi o perdite artificiali di massa. Il programma allegato
contiene parametri, formule, risultati e verifiche. Nessun
parametro \`e stato adattato a immagini o misure di Armando.

\subsection{Alternanze a tempi totali fissi}
Dividendo in $N$ parti le stesse esposizioni complessive si ottiene
\begin{equation}\label{eq:cycles}
 v_N=\left(e^{K_TT_T/N}e^{K_DT_D/N}\right)^Nv_0.
\end{equation}
Non si confrontano quindi quattro cicli lunghi con un solo ciclo
breve. Per $N\to\infty$,
\begin{equation}\label{eq:trotter}
 v_N\longrightarrow e^{K_TT_T+K_DT_D}v_0.
\end{equation}
Il modello mantiene memoria dello stato mobile fra le fasi.
Un modello che lo azzeri a ogni cambio, anche infinitamente rapido,
avrebbe un limite diverso e non sarebbe equivalente.

La Figura~\ref{fig:cycles} confronta tre scenari adimensionali,
con $T_D=T_T=1$, $v_0=(1,0,0)^{\mathsf T}$ e i tassi della
Tabella~\ref{tab:rates}. Il tempo unitario \`e convenzionale:
non significa un minuto o trenta minuti reali.

\begin{table}[htbp]
\caption{Tassi illustrativi $(a,b,e)$ usati nelle simulazioni.
Non sono propriet\`a misurate di DMF o THF.}
\label{tab:rates}\centering\small
\begin{tabular}{@{}lcc@{}}\toprule
Scenario & Trattamento D & Trattamento T\\\midrule
Complementare &(2; 4; 0{,}02)&(0{,}02; 0{,}1; 4)\\
Preparazione poi estrazione &(2; 0; 0)&(0; 0; 2)\\
Tassi uguali &(0{,}4; 0{,}5; 0{,}3)&(0{,}4; 0{,}5; 0{,}3)\\\bottomrule
\end{tabular}
\end{table}

Nel primo scenario il residuo scende da 0,6680 con una sequenza
a 0,4478 con otto alternanze; a venti risale leggermente a 0,4519.
Nel secondo scenario frammentare la preparazione peggiora la
rimozione. Nel terzo non cambia nulla, come deve accadere con
generatori identici. Questi numeri dimostrano possibilit\`a e limiti
del modello, \emph{non} riduzioni previste per il campione.

Un controesempio \`e anche analitico. Se D mobilizza senza ritorno,
T estrae senza ritorno, e $aT_D=kT_T=w$, la frazione rimossa \`e
\begin{equation}\label{eq:counterexample}
 R_N=1-e^{-w}\left[1+N\bigl(1-e^{-w/N}\bigr)\right].
\end{equation}
Essa diminuisce all'aumentare di $N$: preparare tutta la frazione
prima di estrarla \`e qui pi\`u efficace. Cercare un beneficio dei
cicli senza includere questo caso negativo produrrebbe una teoria
costruita per dare ragione alla proposta.

\begin{figure}[htbp]
\centering
\begin{tikzpicture}
\begin{axis}[width=0.94\textwidth,height=6.0cm,xlabel={Numero di alternanze $N$},ylabel={Frazione residua $L+M$},xmin=1,xmax=20,ymin=0,ymax=1,grid=major,legend columns=3,legend style={font=\footnotesize,at={(0.5,1.04)},anchor=south,draw=none},tick label style={font=\small},label style={font=\small}]
\addplot[thick,black,mark=*] coordinates {(1,0.6680496) (2,0.5248837) (3,0.4766819) (4,0.4592498) (5,0.452253) (6,0.4493265) (7,0.4481709) (8,0.4478442) (9,0.4479215) (10,0.4481971) (11,0.4485664) (12,0.4489743) (13,0.4493913) (14,0.4498015) (15,0.4501966) (16,0.4505728) (17,0.4509283) (18,0.451263) (19,0.4515773) (20,0.4518724)};
\addlegendentry{Complementari}
\addplot[thick,blue!75!black,mark=*] coordinates {(1,0.2523549) (2,0.3064317) (3,0.3328908) (4,0.3483364) (5,0.3584219) (6,0.3655152) (7,0.3707725) (8,0.3748238) (9,0.3780406) (10,0.3806565) (11,0.3828253) (12,0.3846526) (13,0.3862129) (14,0.3875609) (15,0.3887371) (16,0.3897723) (17,0.3906905) (18,0.3915104) (19,0.392247) (20,0.3929124)};
\addlegendentry{Preparazione poi estrazione}
\addplot[thick,red!75!black,mark=*] coordinates {(1,0.8798123) (2,0.8798123) (3,0.8798123) (4,0.8798123) (5,0.8798123) (6,0.8798123) (7,0.8798123) (8,0.8798123) (9,0.8798123) (10,0.8798123) (11,0.8798123) (12,0.8798123) (13,0.8798123) (14,0.8798123) (15,0.8798123) (16,0.8798123) (17,0.8798123) (18,0.8798123) (19,0.8798123) (20,0.8798123)};
\addlegendentry{Identici}
\end{axis}
\end{tikzpicture}
\caption{Modello a stati conformazionali. Ogni curva mantiene invariati il tempo totale nel liquido D e quello nel liquido T. Alternare pi\`u spesso pu\`o aiutare o peggiorare: dipende dai tassi, qui illustrativi. D e T sono stati di trattamento ipotetici, non misure di DMF e THF.}
\label{fig:cycles}
\end{figure}

\section{Il solvente deve anche portare via il materiale}
\label{sec:transport}

\subsection{Bilancio fra superficie, strato vicino e bagno}
Sia $\Gamma$ la massa accessibile per area, $I$ una frazione
schermata, $c$ la concentrazione nel bagno, $c_s$ quella prossima
alla superficie. Con $A$ area, $V$ volume, $Q$ portata di ricambio,
poniamo
\begin{align}
 J&=k_d\Gamma-k_a c_s=k_m(c_s-c),\label{eq:flux}\\
 J&=\lambda\Gamma-bc,\qquad
 \lambda=\frac{k_d}{1+k_a/k_m},\quad
 b=\frac{k_a}{1+k_a/k_m},\label{eq:effective}\\
 \dot I&=-k_iI,\quad
 \dot\Gamma=k_iI-J,\quad
 V\dot c=AJ-Qc.\label{eq:balance}
\end{align}
Qui $k_d,k_i$ hanno unit\`a s$^{-1}$; $k_a,k_m$ m/s;
$\Gamma,I$ kg/m$^2$; $c$ kg/m$^3$. Il fluido in entrata \`e
assunto privo del polimero considerato. Si verifica
\begin{equation}
 \frac{d}{dt}\{A(\Gamma+I)+Vc\}=-Qc.
\end{equation}
Il modello assume regime diluito, risposta lineare e parametri
costanti in ciascuna fase. Non rappresenta in dettaglio la forma
delle particelle n\'e una rete reticolata.

Con $h=V/A$, $\tau=\lambda t$, $x=\Gamma/\Gamma_0$,
$y=hc/\Gamma_0$, $z=I/\Gamma_0$,
\begin{equation}\label{eq:dimensionless}
 x'=-x+\beta y+\varepsilon z,\quad
 y'=x-(\beta+r)y,\quad z'=-\varepsilon z,
\end{equation}
dove $\beta=b/(h\lambda)$, $r=Q/(V\lambda)$ e
$\varepsilon=k_i/\lambda$. $\beta$ misura la tendenza a
ripartirsi tra superficie e bagno: \textbf{non \`e la saturazione
del solvente rispetto alla solubilit\`a del PMMA massivo}.

\subsection{Un limite che impedisce di promettere troppo al ricambio}
Da $c,I\ge0$ segue
\begin{equation}\label{eq:bound}
 \Gamma(t)\ge\Gamma_0e^{-\lambda t},\qquad \lambda\le k_d.
\end{equation}
Anche un bagno perfettamente pulito non supera il distacco
intrinseco nel modello. Un trasporto migliore aiuta quando limita
il riadsorbimento; pu\`o essere quasi inutile se domina una
barriera di distacco o una frazione inaccessibile.

Per un bagno senza ricambio, inizialmente pulito e senza $I$,
\begin{equation}\label{eq:static}
 x(\tau)=\frac{\beta+e^{-(1+\beta)\tau}}{1+\beta}.
\end{equation}
Compare un plateau pur essendo tutto reversibile: osservare un
plateau non prova una rete insolubile. La
Figura~\ref{fig:transport} mostra come cambia l'andamento con $r$.

\begin{figure}[htbp]
\centering
\begin{tikzpicture}
\begin{axis}[width=0.94\textwidth,height=6.0cm,xlabel={Tempo adimensionale $\tau=\lambda t$},ylabel={Frazione ancora aderente $x$},xmin=0,xmax=8,ymin=0,ymax=1,grid=major,legend style={font=\footnotesize,at={(0.98,0.98)},anchor=north east},tick label style={font=\small},label style={font=\small}]
\addplot[thick,black] coordinates {(0,1) (0.1,0.9093654) (0.2,0.83516) (0.3,0.7744058) (0.4,0.7246645) (0.5,0.6839397) (0.6,0.6505971) (0.7,0.6232985) (0.8,0.6009483) (0.9,0.5826494) (1,0.5676676) (1.1,0.5554016) (1.2,0.545359) (1.3,0.5371368) (1.4,0.530405) (1.5,0.5248935) (1.6,0.5203811) (1.7,0.5166866) (1.8,0.5136619) (1.9,0.5111854) (2,0.5091578) (2.1,0.5074978) (2.2,0.5061387) (2.3,0.5050259) (2.4,0.5041149) (2.5,0.503369) (2.6,0.5027583) (2.7,0.5022583) (2.8,0.5018489) (2.9,0.5015138) (3,0.5012394) (3.1,0.5010147) (3.2,0.5008308) (3.3,0.5006802) (3.4,0.5005569) (3.5,0.5004559) (3.6,0.5003733) (3.7,0.5003056) (3.8,0.5002502) (3.9,0.5002049) (4,0.5001677) (4.1,0.5001373) (4.2,0.5001124) (4.3,0.5000921) (4.4,0.5000754) (4.5,0.5000617) (4.6,0.5000505) (4.7,0.5000414) (4.8,0.5000339) (4.9,0.5000277) (5,0.5000227) (5.1,0.5000186) (5.2,0.5000152) (5.3,0.5000125) (5.4,0.5000102) (5.5,0.5000084) (5.6,0.5000068) (5.7,0.5000056) (5.8,0.5000046) (5.9,0.5000038) (6,0.5000031) (6.1,0.5000025) (6.2,0.5000021) (6.3,0.5000017) (6.4,0.5000014) (6.5,0.5000011) (6.6,0.5000009) (6.7,0.5000008) (6.8,0.5000006) (6.9,0.5000005) (7,0.5000004) (7.1,0.5000003) (7.2,0.5000003) (7.3,0.5000002) (7.4,0.5000002) (7.5,0.5000002) (7.6,0.5000001) (7.7,0.5000001) (7.8,0.5000001) (7.9,0.5000001) (8,0.5000001)};
\addlegendentry{$r=0$}
\addplot[thick,blue!75!black] coordinates {(0,1) (0.1,0.9093353) (0.2,0.834943) (0.3,0.773743) (0.4,0.7232401) (0.5,0.6814123) (0.6,0.6466215) (0.7,0.6175407) (0.8,0.5930942) (0.9,0.5724109) (1,0.5547847) (1.1,0.5396435) (1.2,0.5265238) (1.3,0.5150498) (1.4,0.5049173) (1.5,0.4958794) (1.6,0.4877362) (1.7,0.4803257) (1.8,0.4735164) (1.9,0.4672019) (2,0.4612957) (2.1,0.4557279) (2.2,0.4504416) (2.3,0.4453907) (2.4,0.4405378) (2.5,0.4358525) (2.6,0.4313101) (2.7,0.4268906) (2.8,0.4225776) (2.9,0.4183577) (3,0.4142203) (3.1,0.4101564) (3.2,0.4061589) (3.3,0.4022217) (3.4,0.3983401) (3.5,0.39451) (3.6,0.3907281) (3.7,0.3869917) (3.8,0.3832986) (3.9,0.3796468) (4,0.3760346) (4.1,0.3724608) (4.2,0.3689243) (4.3,0.3654239) (4.4,0.3619588) (4.5,0.3585283) (4.6,0.3551317) (4.7,0.3517685) (4.8,0.348438) (4.9,0.3451397) (5,0.3418733) (5.1,0.3386383) (5.2,0.3354343) (5.3,0.3322609) (5.4,0.3291178) (5.5,0.3260046) (5.6,0.3229211) (5.7,0.3198668) (5.8,0.3168416) (5.9,0.3138451) (6,0.3108769) (6.1,0.3079369) (6.2,0.3050248) (6.3,0.3021402) (6.4,0.299283) (6.5,0.2964528) (6.6,0.2936493) (6.7,0.2908724) (6.8,0.2881218) (6.9,0.2853972) (7,0.2826984) (7.1,0.2800251) (7.2,0.2773771) (7.3,0.2747541) (7.4,0.272156) (7.5,0.2695824) (7.6,0.2670331) (7.7,0.264508) (7.8,0.2620067) (7.9,0.2595291) (8,0.2570749)};
\addlegendentry{$r=0.2$}
\addplot[thick,red!75!black] coordinates {(0,1) (0.1,0.909078) (0.2,0.8331708) (0.3,0.7685775) (0.4,0.7126311) (0.5,0.6634017) (0.6,0.6194848) (0.7,0.5798518) (0.8,0.5437427) (0.9,0.5105898) (1,0.4799642) (1.1,0.4515367) (1.2,0.4250503) (1.3,0.4003011) (1.4,0.3771236) (1.5,0.3553811) (1.6,0.3349583) (1.7,0.3157562) (1.8,0.2976884) (1.9,0.2806782) (2,0.2646569) (2.1,0.2495622) (2.2,0.2353369) (2.3,0.2219286) (2.4,0.2092886) (2.5,0.1973715) (2.6,0.1861352) (2.7,0.1755401) (2.8,0.1655492) (2.9,0.1561277) (3,0.147243) (3.1,0.1388643) (3.2,0.1309626) (3.3,0.1235107) (3.4,0.116483) (3.5,0.1098553) (3.6,0.1036048) (3.7,0.09770993) (3.8,0.09215052) (3.9,0.08690745) (4,0.08196271) (4.1,0.07729933) (4.2,0.07290128) (4.3,0.06875348) (4.4,0.06484167) (4.5,0.06115243) (4.6,0.0576731) (4.7,0.05439173) (4.8,0.05129706) (4.9,0.04837846) (5,0.04562592) (5.1,0.04302999) (5.2,0.04058176) (5.3,0.03827282) (5.4,0.03609525) (5.5,0.03404158) (5.6,0.03210475) (5.7,0.03027812) (5.8,0.02855542) (5.9,0.02693073) (6,0.02539848) (6.1,0.02395341) (6.2,0.02259056) (6.3,0.02130525) (6.4,0.02009307) (6.5,0.01894985) (6.6,0.01787168) (6.7,0.01685486) (6.8,0.01589589) (6.9,0.01499147) (7,0.01413852) (7.1,0.0133341) (7.2,0.01257544) (7.3,0.01185995) (7.4,0.01118517) (7.5,0.01054878) (7.6,0.009948593) (7.7,0.009382559) (7.8,0.008848729) (7.9,0.008345273) (8,0.007870461)};
\addlegendentry{$r=2$}
\addplot[thick,green!45!black] coordinates {(0,1) (0.1,0.9074075) (0.2,0.8249234) (0.3,0.7501247) (0.4,0.6821311) (0.5,0.6203034) (0.6,0.5640801) (0.7,0.5129527) (0.8,0.4664595) (0.9,0.4241804) (1,0.3857334) (1.1,0.3507711) (1.2,0.3189778) (1.3,0.2900662) (1.4,0.2637751) (1.5,0.2398669) (1.6,0.2181258) (1.7,0.1983552) (1.8,0.1803766) (1.9,0.1640276) (2,0.1491604) (2.1,0.1356407) (2.2,0.1233465) (2.3,0.1121665) (2.4,0.1019999) (2.5,0.09275484) (2.6,0.08434769) (2.7,0.07670255) (2.8,0.06975036) (2.9,0.0634283) (3,0.05767926) (3.1,0.05245131) (3.2,0.0476972) (3.3,0.04337401) (3.4,0.03944266) (3.5,0.03586764) (3.6,0.03261665) (3.7,0.02966033) (3.8,0.02697197) (3.9,0.02452727) (4,0.02230416) (4.1,0.02028255) (4.2,0.01844417) (4.3,0.01677242) (4.4,0.0152522) (4.5,0.01386976) (4.6,0.01261263) (4.7,0.01146944) (4.8,0.01042987) (4.9,0.009484523) (5,0.008624861) (5.1,0.007843118) (5.2,0.007132231) (5.3,0.006485777) (5.4,0.005897916) (5.5,0.005363339) (5.6,0.004877215) (5.7,0.004435152) (5.8,0.004033157) (5.9,0.003667598) (6,0.003335173) (6.1,0.003032878) (6.2,0.002757983) (6.3,0.002508004) (6.4,0.002280683) (6.5,0.002073965) (6.6,0.001885984) (6.7,0.001715042) (6.8,0.001559593) (6.9,0.001418234) (7,0.001289688) (7.1,0.001172793) (7.2,0.001066493) (7.3,0.0009698274) (7.4,0.0008819238) (7.5,0.0008019877) (7.6,0.0007292968) (7.7,0.0006631945) (7.8,0.0006030837) (7.9,0.0005484211) (8,0.0004987131)};
\addlegendentry{$r=20$}
\addplot[black,dashed,domain=0:8,samples=80]{exp(-x)};
\addlegendentry{Estrazione ideale}
\end{axis}
\end{tikzpicture}
\caption{Modello di trasporto: $\beta=1$, con diversi ricambi $r$. Parametri scelti, non stimati sul campione. Il ricambio riduce il riadsorbimento ma non supera la velocit\`a intrinseca di distacco.}
\label{fig:transport}
\end{figure}

\subsection{Ricambi a volume e tempo totali invariati}
Nel caso senza frazione schermata, $I_0=0$, se si dividono
$V_{\rm tot}$ e $t_{\rm tot}$ in $N$ bagni freschi
uguali, mantenendo $\lambda$ invariato per ipotesi e ponendo
$B=Ak_a/(V_{\rm tot}k_d)$, il residuo \`e
\begin{equation}\label{eq:resource}
 x_N=\left[\frac{NB+e^{-(1+NB)\tau_{\rm tot}/N}}{1+NB}\right]^N,
 \quad \tau_{\rm tot}=\lambda t_{\rm tot}.
\end{equation}
Per $B>0$,
\begin{equation}\label{eq:resource-limit}
 \lim_{N\to\infty}x_N=
 \exp\!\left[-\frac{1-e^{-B\tau_{\rm tot}}}{B}\right].
\end{equation}
Per $B=0$ il risultato \`e $e^{-\tau_{\rm tot}}$ per ogni $N$.
Ad esempio, con $B=1$ e $\tau_{\rm tot}=5$, passare da uno a
quattro bagni porta il residuo da 0,5000 a 0,4104; il limite
di infiniti ricambi \`e 0,3704. Sono valori del modello.
La formula non giustifica bagni arbitrariamente piccoli in
laboratorio: cambiare geometria pu\`o cambiare $k_m$ e $\lambda$.

\subsection{Un parametro fisico verificato, con un uso delimitato}
Arai e collaboratori \cite{Arai1996} misurano
$D=2{,}93\times10^{-11}\,\mathrm{m^2/s}$ per PMMA atattico
di $M_w=9{,}52\times10^5$ in acetone a 25\,\degC.
\`E la diffusione di catene gi\`a sciolte. Per distanze
ipotetiche $\delta$ di 10, 100 e 1000\,\um,
$\delta^2/D$ vale circa 3,4 secondi, 5,7 minuti e 9,5 ore.
Questi ordini di grandezza mostrano il peso della geometria.
Non determinano la diffusione in DMF/THF, il tempo di ingresso
nel residuo, $k_d$, o uno spessore $\delta$ dai soli 500 rpm.

\section{Una spiegazione dell'asciugatura che si pu\`o smentire}

Se il liquido trattenuto sul campione all'uscita ha concentrazione
massica $c_f$ e volume $V_f$, la massa che pu\`o depositare \`e
limitata da
\begin{equation}\label{eq:mass-bound}
 M_{\rm dep}\le c_fV_f,\qquad
 \Gamma_{\rm dep}\le c_fh_f,\quad h_f=V_f/A.
\end{equation}
Per uno spessore equivalente e densit\`a $\rho$,
\begin{equation}\label{eq:units}
 t_{\rm eq}[\mathrm{nm}]\le
 \frac{0{,}001\;c_f[\mathrm{mg/L}]\;h_f[\mu\mathrm m]}
 {\rho[\mathrm{g/cm^3}]}.
\end{equation}
L'esempio $c_f=1$ mg/L, $h_f=10$\,\um\ e densit\`a
assunta 1,18 g/cm$^3$ d\`a al massimo 0,0085 nm medi.
Per spiegare 0,5 nm medi, con quello stesso film, servirebbero
almeno 59 mg/L. Non sono misure di questo campione.

Per usare il limite occorre delimitare la concentrazione
\emph{locale} del liquido trascinato: la media di un bagno grande
pu\`o sottostimarla. Il confronto va fatto su massa e area
complessive; l'evaporazione pu\`o concentrare il materiale in
isole alte. Il limite esclude soltanto il deposito proveniente
dall'ultimo film liquido, non il riadsorbimento avvenuto durante
tutto il bagno.

Questa misura decide se valga la pena ottimizzare l'asciugatura.
Se perfino il massimo materiale trascinabile \`e insufficiente,
``\`e tutto un deposito finale'' non spiega il residuo.
Se \`e sufficiente, il meccanismo rimane possibile ma non provato.

\section{Il programma sperimentale proposto}
\label{sec:experiment}

\subsection{Prima prova: separare i due passaggi IPA}
Si usano campioni gemelli dello stesso trasferimento e la stessa
temperatura del procedimento di Armando, da registrare. Ogni
numero della Tabella~\ref{tab:factorial} \`e un bagno distinto;
i bagni sostitutivi sono freschi. Volume, geometria, agitazione,
intervallo di trasferimento e getto di azoto devono essere
uguali. Non si lasciano asciugare i campioni fra due bagni.
R riproduce la sequenza riferita con freschezza dei bagni
standardizzata; non sappiamo se questo dettaglio coincida con
le prove gi\`a eseguite da Armando.

\begin{table}[htbp]
\caption{Quattro percorsi da 76 minuti. ``5+3'' significa due
bagni freschi distinti. Nessun risultato \`e ancora acquisito.}
\label{tab:factorial}\centering\small
\begin{tabular}{@{}lcccc@{}}\toprule
Gruppo & Primo & Intermedio & Secondo & Finale\\\midrule
R, riferimento &DMF 30&IPA 5+3&THF 30&IPA 5+3\\
I, cambia intermedio &DMF 30&DMF 5+3&THF 30&IPA 5+3\\
F, cambia finale &DMF 30&IPA 5+3&THF 30&THF 5+3\\
IF, cambia entrambi &DMF 30&DMF 5+3&THF 30&THF 5+3\\\bottomrule
\end{tabular}
\end{table}

Le varianti I e IF costituiscono il primo tentativo concreto di
evitare un eventuale riblocco prima dell'estrazione in THF.
Il confronto I contro R e IF contro F misura la sostituzione
intermedia nei due contesti. F contro R e IF contro I studia
il passaggio finale, che cambia anche evaporazione e forze
capillari. Nessuna di queste differenze va chiamata automaticamente
``precipitazione''.

Si raccoglie separatamente il liquido dopo DMF, dopo ciascun IPA
e dopo THF. Su aliquote dell'eluato reale, fuori dal campione,
si verifica l'eventuale comparsa di particelle durante l'aggiunta
di IPA, usando bianchi e PMMA/PPC dello stesso lotto. Si cerca
anche se il materiale torna solubile reintroducendo il liquido
di partenza. L'esito pu\`o motivare un percorso di scambio
graduale, ma non fissa un rapporto ottimale senza misure.

Un esito negativo della prova di torbidit\`a non esclude
aggregati molto piccoli o nucleazione selettiva sulla superficie.
Un bianco di solo \SiO\ non riproduce tutte le propriet\`a
del grafene: serve anche, se disponibile, un testimone grafenico
pulito sottoposto al medesimo eluato.

\subsection{Seconda prova: sfruttare un'eventuale mobilizzazione reversibile}
Se l'IPA intermedio \`e utile, oppure se si osserva riorganizzazione
seguita da esportazione nel THF, si prova la
Tabella~\ref{tab:cycles}. L'ipotesi \`e estrarre presto la
frazione mobile prima che ritorni aderente.

\begin{table}[htbp]
\caption{Confronto fra blocchi e alternanze, entrambi da 76 minuti:
30 minuti DMF, 16 IPA, 30 THF. Stesso numero di bagni freschi,
stesso liquido finale. Il riferimento R va misurato separatamente.}
\label{tab:cycles}\centering\small
\begin{tabular}{@{}p{0.14\textwidth}p{0.77\textwidth}@{}}\toprule
Percorso & Sequenza\\\midrule
Blocchi & DMF 30; quattro bagni IPA da 4 minuti;
quattro bagni THF da 7,5 minuti; azoto.\\[4pt]
Alternato & DMF 30; quattro ripetizioni di IPA 4 minuti
seguito da THF 7,5 minuti; azoto.\\\bottomrule
\end{tabular}
\end{table}

Non \`e un invito a moltiplicare indiscriminatamente i bagni.
Le due varianti hanno nove bagni ciascuna, incluso il DMF
iniziale, e terminano entrambe nel THF. Il confronto tiene
fissi il tempo complessivo in ogni liquido e il numero di
manipolazioni; l'ordine resta la variabile. Per ciascuna coppia
si usano volumi, recipienti e modalit\`a di ricambio identici.

Il materiale esportato nei singoli bagni \`e il primo esito
da confrontare, insieme alla superficie conservata. Se
l'alternanza migliora solo l'immagine finale ma non le altre
misure, non \`e ancora l'amplificazione del meccanismo proposto.
La ricerca di una durata ottimale dei cicli ha senso soltanto
dopo un segnale replicato a tempi totali uguali.

\subsection{Una previsione geometrica aggiuntiva}
Se domina una rimozione dal bordo di isole isolate e di altezza
costante, con velocit\`a $v$, allora
\begin{equation}\label{eq:edge}
 r(t)=\max(r_0-vt,0),\qquad
 \frac{M(t)}{M_0}=\left[\max\left(1-\frac{vt}{r_0},0\right)\right]^2.
\end{equation}
Il tempo di sparizione cresce come $r_0$. Una rimozione
volumetrica di primo ordine prevede invece lo stesso tempo
di riduzione frazionale per tutte le dimensioni. Si seguono
quindi anche raggio, altezza e volume delle stesse isole.
Una dipendenza dal raggio non dimostra da sola il meccanismo
di bordo: diffusione, eterogeneit\`a e coalescenza vanno distinti.

\section{Le misure che decidono la riuscita}

\subsection{Identit\`a: un controllo breve prima di inseguire altri liquidi}
Si confrontano residuo persistente, eluati e riferimenti PMMA,
PPC e loro sovrapposizione, trattati con la stessa storia.
Se il riferimento libero si dissolve ma il residuo sul grafene
resta, diventa plausibile una barriera interfaciale. Se persiste
anche il riferimento privo di grafene, occorre verificare
alterazione, reticolazione o materiale diverso. I 90\,\degC\
per due minuti, da soli, non provano carbonizzazione.

Un singolo picco carbonilico XPS non separa PMMA e PPC.
Si usano firme complete e, se la quantit\`a lo consente,
spettroscopia infrarossa locale o frammenti ToF-SIMS confrontati
con riferimenti. Le analisi distruttive richiedono gemelli
dedicati. I materiali ausiliari, come PDMS o adesivi, entrano
nei controlli soltanto se sono effettivamente usati. Il segnale
globale di silicio su \SiO\ non identifica un residuo siliconico.

Per studiare un'eventuale modifica chimica si analizzano anche
riferimenti esposti ai bagni. La perdita della firma dell'estere
pu\`o indicare trasformazione senza asportazione. La cromatografia
dimensionale del materiale estratto, se misurabile, riguarda la
frazione solubile e non esclude una rete insolubile rimasta sul
campione. Nessun singolo strumento chiude da solo il bilancio.

\subsection{Morfologia, integrit\`a e limite di rivelabilit\`a}
Si misurano copertura, volume apparente per area e firma chimica,
non soltanto rugosit\`a. Una pellicola uniforme pu\`o essere
liscia e ancora presente. Le stesse aree vengono registrate
prima e dopo, insieme ad aree casuali e campioni non
prescansionati. La punta AFM non deve diventare l'agente di
pulizia; in liquido si confrontano campioni nello stesso mezzo.

La rimozione va calcolata sull'intersezione delle regioni dove
il grafene \`e presente prima e dopo. Le regioni asportate
sono danno, non superficie pulita. Mappa Raman, fori,
delaminazioni e substrati testimone controllano l'integrit\`a.
Mobilit\`a, resistenza di contatto e drogaggio restano esiti
separati dalla quantit\`a di polimero.

Come primo obiettivo progettuale si pu\`o fissare una riduzione
del 90\% di copertura e volume residuo, con diminuzione
chimica coerente e almeno 99\% dell'area di grafene conservata.
Non sono risultati ottenuti n\'e standard universali. Per la
scommessa occorre concordare prima se basta tale miglioramento
oppure se si richiede residuo non rilevato a un limite dichiarato.
Quest'ultimo richiede calibrazione della sensibilit\`a e non
equivale ad assenza assoluta di ogni molecola.

Una prima esplorazione pu\`o usare un trasferimento suddiviso
fra condizioni; la conferma deve coinvolgere trasferimenti
indipendenti, inizialmente almeno tre e poi un numero adeguato
alla variabilit\`a osservata. I pixel e le immagini dello
stesso campione non sono repliche indipendenti. Si riportano
tutte le condizioni, incluse quelle che non migliorano.


\section{Calcoli molecolari eseguiti e ipotesi scartate}\label{sec:molecular}

\subsection{Il bersaglio locale non determina il distacco}

Il PMMA presenta gruppi estere; PPC presenta gruppi carbonato. Abbiamo
calcolato interazioni di HFIP con acetato di metile, pivalato di metile
e carbonato di dimetile come analoghi locali, oltre a coppie di solventi.
Queste molecole non rappresentano la catena 950K, le sue molteplici
ancore o la storia del deposito. Sono controlli di un'ipotesi locale.

Sono state ottimizzate 12 molecole e 13 coppie, con tre geometrie iniziali
per coppia, usando GFN2-xTB. Tutte le 39 geometrie delle coppie sono state
raffinate fino a convergenza con forza massima sotto
0,008\,eV/\AA. Non rivendichiamo un minimo globale.
Le energie di interazione a geometria fissata sono
\begin{equation}
 E_{\rm int}=E_{AB}(\mathbf R_A,\mathbf R_B)
             -E_A(\mathbf R_A)-E_B(\mathbf R_B),
\end{equation}
separate dall'energia di legame che comprende la deformazione dei frammenti.

Abbiamo poi eseguito 50 calcoli autoconsistenti PBE0-D3(BJ), tutti
convergenti, per sette coppie con base def2-SVP e tre controlli con
def2-TZVP. La Tabella~\ref{tab:quantum} corregge l'errore dovuto
all'uso della base del partner tramite il metodo \emph{counterpoise}.
I frammenti rimangono nella geometria del complesso ottenuta con GFN2.
Sono energie elettroniche nel vuoto, non energie libere nel liquido.

\begin{table}[htbp]\centering
\caption{Interazione con acetato di metile a geometria fissata,
PBE0-D3(BJ); kJ/mol. La correzione di base \`e positiva e riduce
l'apparente attrazione. Le cifre esprimono il calcolo numerico,
non un'accuratezza fisica al centesimo.}\label{tab:quantum}
\small
\begin{tabular}{lcrr}\toprule
Partner & Base & $E_{\rm int}$ corretto & Correzione di base\\\midrule
HFIP & SVP & -45.28 & 20.17 \\
HFIP & TZVP & -44.34 & 2.86 \\
Acido acetico & SVP & -48.69 & 15.00 \\
Acido acetico & TZVP & -47.53 & 2.00 \\
IPA & SVP & -28.92 & 14.17 \\
IPA & TZVP & -28.01 & 2.52 \\
\bottomrule\end{tabular}
\end{table}

\textbf{Il confronto produce una confutazione utile: anche l'acido
acetico interagisce almeno altrettanto fortemente con questo analogo,
pur essendo gi\`a fallito sul materiale riferito.} Per queste geometrie
la differenza fra SVP e TZVP dopo correzione \`e circa 1\,kJ/mol;
non dimostra per\`o convergenza rispetto a geometrie, funzionale,
entropia, liquido o superficie. Il solo legame con l'estere non
discrimina un trattamento efficace.

La letteratura documenta PMMA di massa 99.200\,g/mol disciolto in HFIP
a temperatura ambiente \cite{WangHFIP2020}. Dimostra la solvatazione
di quel materiale libero. La selezione di HFIP perch\'e dissolva anche
metilcellulosa non implica che gli altri solventi non dissolvano PMMA.
Non dimostra rimozione del PMMA 950K aderente al grafene o del PPC.

Un eventuale primo confronto operativo usa HFIP puro in due bagni
freschi di 10 minuti a temperatura ambiente, con il secondo destinato
a esportare il materiale estratto, inizialmente evitando la mescolanza
con IPA. \textbf{Quei tempi sono una proposta pilota, non condizioni
validate o derivate dai calcoli.} Occorrono un controllo del materiale
libero dello stesso lotto, del substrato e della contaminazione
fluorurata dopo il trattamento, insieme alle misure della Sezione~9.
La sequenza abituale conserva la funzione di controllo operativo.

\subsection{Perch\'e serve il liquido e una selettivit\`a interfaciale}

Siano $P$ il polimero, $G$ il grafene, $O$ l'ossido e $L$ il liquido.
Nel limite reversibile e non reattivo i lavori di separazione sono
\begin{align}
 W_{PG}^{L}&=\gamma_{PL}+\gamma_{GL}-\gamma_{PG},\\
 W_{GO}^{L}&=\gamma_{GL}+\gamma_{OL}-\gamma_{GO}.
\end{align}
La differenza cancella $\gamma_{GL}$: una maggiore affinit\`a del
solvente per il grafene non garantisce selettivit\`a fra rimozione
del polimero e distacco del foglio dall'ossido. Questi lavori non
sono, da soli, barriere cinetiche o energie di frattura del processo
di prelievo, che dipende anche da geometria, temperatura e velocit\`a.

Il medesimo problema appare a scala molecolare. Se il carbonile \`e
solvatato nello stesso modo sia nel polimero aderente sia nel polimero
liberato, quel contributo si cancella nel costo del distacco.
Bisogna calcolare la differenza fra i due stati, includendo liquido,
contatti, infiltrazione, conformazioni e riaggancio. Moltiplicare
un'energia di coppia per circa 9.500 unit\`a non fornisce una barriera
corretta: la catena pu\`o staccarsi progressivamente.

\subsection{Un controllo del liquido che limita le conclusioni}

Abbiamo trascritto i parametri HFIP pubblicati nelle Tabelle I--V
del supplemento di Marchelli et al. \cite{Marchelli2022}, convertendo
esplicitamente unit\`a e prefattori armonici. Il fattore 1--4, non
specificato nel testo consultato, \`e assunto pari a 0,5 e deve essere
oggetto di sensibilit\`a. Sono state simulate 125 molecole:
50\,ps a volume fissato e 1\,ns a pressione fissata a 298,15\,K.
L'analisi usa 200--1000\,ps; una sola realizzazione e questo intervallo
non bastano a certificare convergenza della risposta dielettrica.

\begin{table}[htbp]\centering
\caption{Controllo preliminare del modello HFIP. I riferimenti
sperimentali sono riportati da Casoria et al. \cite{Casoria2024}.
La deviazione della densit\`a \`e uno scarto temporale, non un
intervallo di confidenza di repliche indipendenti.}
\begin{tabular}{lcc}\toprule
Quantit\`a & Simulazione pilota & Riferimento sperimentale\\\midrule
Densit\`a, g/cm$^3$ & 1,572 ($s=0,022$) & 1,607\\
Permittivit\`a relativa & 1,635 & 16,7\\\bottomrule
\end{tabular}
\end{table}

L'espressione delle fluttuazioni del dipolo \`e stata verificata
indipendentemente in unit\`a SI:
\begin{equation}
 \epsilon_r=1+\frac{\langle\mathbf M^2\rangle
                         -\langle\mathbf M\rangle^2}
                       {3\epsilon_0 k_B T\langle V\rangle},
\end{equation}
con condizioni elettrostatiche conduttrici al contorno e cariche fisse.
Il dipolo dell'ultimo fotogramma, ricostruito separatamente dalle
coordinate e dai parametri pubblicati, concorda entro
0,0005\,Debye con il programma. Il contributo elettronico non \`e
incluso da un modello a cariche fisse; non basta a giustificare
automaticamente il divario osservato.

\textbf{La vicinanza della densit\`a non qualifica il modello per
una classifica quantitativa del distacco.} Vanno controllate durata,
repliche, convenzioni 1--4, conformazione del solvente e polarizzazione.
Una conversione errata del volume nelle prime esecuzioni \`e stata
rilevata mediante controllo fra motori e derivazione indipendente
in SI: quei tentativi sono archiviati come invalidi ed esclusi.

\subsection{Potenziali appresi e catene realistiche: stato della ricerca}

MACE-OFF23-small \cite{MACE2025} \`e stato effettivamente eseguito
sulle 13 coppie e confrontato con i calcoli precedenti. Usa un diverso
livello quantistico di addestramento, quindi le discrepanze non sono
una misura universale dell'errore rispetto a PBE0. Il modello non
comprende silicio: non descrive l'intero supporto Si/\SiO.
La sua esecuzione Metal a precisione singola \`e stata confrontata
con CPU a precisione doppia; la differenza quadratica media delle
forze sul complesso di controllo \`e circa
$1,9\times10^{-6}$\,eV/\AA. Una modifica locale, registrata,
riguarda soltanto il tipo numerico delle energie atomiche diagnostiche
su Metal, senza modificare l'espressione di energia totale e forze.

Sono inoltre state costruite topologie di PMMA a 4 e 8 unit\`a con
ogni termine di legame, angolo e torsione vincolato alla Tabella S2
di Behbahani e Harmandaris \cite{Behbahani2021}. Mancanze avrebbero
prodotto un errore esplicito. La neutralizzazione di 0,0051 elettroni
in eccesso ai due estremi \`e una scelta di chiusura documentata,
ancora da verificare. Quelle catene non sono ancora una simulazione
del distacco e non sostituiscono la massa 950K.

\textbf{Stato: traiettorie interfaciali eseguite su un oligomero;
nessuna energia libera di distacco o pulizia reale ancora qualificata.}
\`E stato eseguito anche MACE-POLAR-1-small \cite{MACEPolar2026},
con interazioni elettrostatiche e 83 elementi. Per gli stessi complessi
con acetato di metile otteniamo $-46{,}19$\,kJ/mol con HFIP,
$-44{,}83$ con acido acetico e $-28{,}27$ con IPA. Il livello di
addestramento \`e diverso da PBE0: la variazione dell'ordine fra
HFIP e acido, di pochi kJ/mol, non produce una previsione di pulizia.
La versione della libreria elettrostatica \`e stata fissata a 0.4.0
per corrispondere all'interfaccia del modello rilasciato.

Un controllo ulteriore applica un campo statico alla stessa molecola
HFIP, a nuclei fissi. PBE0/def2-TZVP produce una polarizzabilit\`a
lungo $z$ di circa $0{,}4242\,e\,\text{\AA}^2/\mathrm V$, coerente fra
ampiezze 0,002 e 0,01\,V/\AA. Nel percorso di esecuzione del
modello appreso otteniamo invece circa $-0{,}0335\,e\,\text{\AA}^2/\mathrm V$
a 0,01\,V/\AA. La stima usa
\begin{equation}
 \alpha_{zz}\simeq -\frac{E(F)+E(-F)-2E(0)}{F^2}.
\end{equation}
Per uno stato fondamentale elettronico stabile, l'energia ottenuta
minimizzando un funzionale con accoppiamento lineare al campo \`e
concava in $F$. Il segno negativo ottenuto dal modello appreso
richiede chiarire convenzioni, implementazione e dominio di validit\`a;
la risposta al campo di questa configurazione non viene promossa a
previsione fisica del liquido. Una buona energia di coppia non
sostituisce questo controllo.


\subsection{Controllo della dimensione del modello e della rotazione di HFIP}

L'anomalia della versione piccola non viene generalizzata all'intera
famiglia. Ripetendo il controllo con MACE-POLAR-1-M, a nuclei fissi,
otteniamo diagonali positive per HFIP, DMF e THF. Per HFIP
$\alpha_{zz}=0{,}4947\,e\,\text{\AA}^2/\mathrm V$, da confrontare
con $0{,}4242$ di PBE0. La traslazione della molecola neutra non cambia
l'energia oltre $10^{-12}$\,eV. Abbiamo verificato nel codice la
convenzione: la molteplicit\`a del singoletto \`e 1, mentre il numero
di elettroni spaiati \`e zero. Anche con le geometrie centrate la
versione S d\`a diagonali negative per tutte e tre le molecole.
Il problema \`e quindi specifico a quella parametrizzazione o architettura,
non una conclusione contro tutti i potenziali appresi.

Abbiamo ruotato soltanto l'idrogeno OH in dodici posizioni, mantenendo
fissi tutti gli altri nuclei e confrontando le energie relative.
La barriera campionata \`e 15,075\,kJ/mol con
PBE0-D3(BJ)/def2-TZVP, 14,145 con MACE-POLAR-1-M e 15,320 con S.
GAFF2 d\`a 17,793\,kJ/mol, ma cambia anche la forma del profilo:
il massimo \`e a $0^\circ$, mentre nei riferimenti quantistici
\`e vicino a $\pm90^\circ$. Il controllo riguarda la molecola isolata;
non fornisce la popolazione dei conformeri nel liquido o la sua
permittivit\`a. Il modello M supera questi controlli molecolari,
senza che ne segua una qualificazione di grafene o \SiO.

Il controllo \`e stato esteso a GAFF2/AM1-BCC per HFIP e DMF:
125 molecole, 500\,ps a pressione fissata dopo 50\,ps iniziali.
Nella parte analizzata la densit\`a \`e circa 1,623 e 0,965\,g/cm$^3$,
mentre le stime della permittivit\`a sono 5,96 e 19,84.
Il riferimento sperimentale DMF a 298,15\,K e 0,1\,MPa \`e
$37,2\pm0,1$ \cite{Hunger2010}. Anche questi controlli hanno
una sola realizzazione breve: non promuovono una previsione
quantitativa del distacco dal buon accordo della sola densit\`a.

\subsection{Traiettorie effettive del polimero sulla superficie}

Sono state eseguite cinque traiettorie a 298,15\,K con quattro unit\`a
PMMA, 65 atomi, e un foglio periodico rigido di 448 atomi di carbonio.
I parametri PMMA e grafene vengono dal supplemento di
Behbahani e Harmandaris \cite{Behbahani2021}; i solventi
da GAFF2/AM1-BCC, con neutralit\`a verificata. La cella misura
$3{,}4032\times3{,}4385\times3{,}8$\,nm. Seguiamo 50\,ps iniziali
a volume fisso e 100\,ps di osservazione per liquido.

\textbf{Il sistema contiene liquido su entrambe le facce del foglio:
non simula ancora grafene appoggiato sull'ossido.} L'assenza
di \SiO\ limita sia la selettivit\`a sia l'affinit\`a quantitativa:
il substrato reale pu\`o modificare il campo e le interazioni attraverso
il monostrato. Il grafene viene tenuto fisso; la sua immobilit\`a
verifica il vincolo numerico, non l'assenza di danni del trattamento.
La catena corta non rappresenta PMMA 950K, un deposito multistrato,
una frazione covalente o un residuo sepolto.

Nessuna di queste traiettorie produce distacco completo spontaneo.
La Tabella~\ref{tab:interface} scompone l'interazione diretta
polimero--grafene, mediata sulla seconda met\`a dei 100\,ps.
\textbf{Questa energia non \`e l'energia libera di adsorbimento.}
Non include l'intero bilancio polimero--liquido e grafene--liquido,
n\'e l'entropia o una dimostrazione di equilibrio.

\begin{table}[htbp]\centering\small
\caption{Energie dirette polimero--grafene nel modello con quattro
unit\`a PMMA, in kJ/mol di oligomero. Una sola realizzazione per
liquido; medie temporali, non una classifica di pulizia.}
\label{tab:interface}
\begin{tabular}{lrrrr}\toprule
Liquido & Metile estere & Nucleo estere & Scheletro/metili & Totale\\\midrule
HFIP & -14,87 & -17,76 & -35,38 & -68,01 \\
DMF & -9,45 & -16,52 & -36,18 & -62,15 \\
THF & -14,04 & -21,47 & -35,86 & -71,37 \\
Acido acetico & -5,94 & -13,00 & -45,48 & -64,42 \\
Anisolo & -1,95 & -7,24 & -38,10 & -47,29 \\
\bottomrule\end{tabular}
\end{table}

In anisolo i gruppi estere perdono molti contatti, mentre lo scheletro
rimane aderente. In HFIP il residuo aderente forma gi\`a circa
2,9 legami idrogeno con i carbonili, contro circa 3,5 nello stato
distaccato ottenuto con un vincolo di separazione. Questa osservazione
del modello \`e coerente con una solvatazione gi\`a presente nello
stato aderente: contare legami idrogeno non basta a stabilire il vantaggio
del distacco. Le distribuzioni conformazionali e la loro convergenza
rimangono da controllare. L'anisolo \`e inoltre il solvente del
PMMA 950K commerciale; ci\`o motiva il controllo, ma la solubilit\`a
del materiale libero non dimostra il recupero del campione.

\subsection{Separazione controllata e motivo del mancato profilo di energia libera}

Per HFIP, DMF e anisolo abbiamo eseguito tre finestre di 125\,ps,
con centri a 0,60, 0,90 e 1,35\,nm dal foglio e molla
$K=2000$\,kJ\,mol$^{-1}$\,nm$^{-2}$. A 1,35\,nm la seconda
met\`a di ciascun campione non ha contatti stretti con il foglio,
e l'interazione diretta residua \`e trascurabile. \`E un distacco
forzato di un oligomero, non una pulizia ottenuta da un bagno.
La forza restituita dal programma coincide con quella ricostruita
da $K(z_0-\langle z\rangle)$.

Le tre distribuzioni di distanza non hanno sovrapposizione osservata
fra finestre adiacenti. Non si possono collegare i loro livelli
di energia libera senza ulteriori campioni. Inoltre, nella finestra
HFIP a 1,35\,nm la forza media passa da circa $+134$ a $-70$
kJ\,mol$^{-1}$\,nm$^{-1}$ fra blocchi iniziale e finale di 25\,ps:
\`e una deriva, non una differenza di equilibrio rispetto a DMF.
\textbf{Non viene quindi calcolato un costo di distacco dai tre punti.}

La campagna pi\`u fitta usa undici finestre di 500\,ps,
con $K=750$ e centri da 0,45 a 1,45\,nm. Le undici finestre
HFIP e le undici DMF sono completate: in HFIP rimane una
disconnessione osservata fra i centri 0,75 e 0,85\,nm. Sono
completate anche cinque finestre inverse HFIP, da 1,05 a 0,65\,nm;
le due storie conservano differenze e lacune di campionamento. Un controllo delle conformazioni rivela inoltre
il limite del solo centro di massa, descritto qui sotto.
Il codice che ricostruisce le distribuzioni \`e stato verificato
su campioni indipendenti di un potenziale armonico noto:
errore quadratico medio del profilo 0,068\,kJ/mol.
Questo qualifica il controllo numerico, non le traiettorie del materiale.
Nessun tempo di pulizia sperimentale viene letto dai percorsi accelerati.

\subsection{Una variabile nascosta osservata: l'orientazione della catena}

Abbiamo confrontato due campioni HFIP di 500\,ps con lo stesso
vincolo $z_0=1{,}05$\,nm e $K=750$, avviati dalla catena aderente
e dalla catena libera. Negli ultimi 250\,ps la distribuzione della
distanza ha sovrapposizione empirica circa 0,355; la distribuzione
congiunta di distanza ed estensione normale non ha sovrapposizione
nelle celle usate, di 0,01 e 0,005\,nm. La catena conserva
orientazioni diverse. Questo \`e un risultato del modello e
delle traiettorie effettive, non una barriera del residuo reale.

Definiamo l'estensione normale con le masse $m_i$:
\begin{equation}
 q=R_{g,z}=\left[\frac{\sum_i m_i(z_i-z_{\rm CM})^2}
 {\sum_i m_i}\right]^{1/2}.
\end{equation}
Per isolare l'effetto a distanze effettivamente uguali consideriamo
le configurazioni con $1{,}00\le z<1{,}04$\,nm.
Le rispettive altezze medie coincidono entro $8\times10^{-5}$\,nm,
mentre cambiano estensione e prossimit\`a al foglio
(Tabella~\ref{tab:orientation}). I punti sono correlati;
le medie non sono stime su repliche indipendenti.

\begin{table}[ht]
\centering\small
\caption{Confronto a pari intervallo di altezza, nel modello
PMMA4/grafene rigido/HFIP. L'energia \`e l'interazione diretta
Lennard--Jones, non un'energia libera di adsorbimento.}
\label{tab:orientation}
\begin{tabular}{lrr}
\toprule
 & Dallo stato aderente & Dallo stato libero\\
\midrule
Configurazioni salvate &52&63\\
$\langle z\rangle$ (nm)&1,02075&1,02067\\
$\langle R_{g,z}\rangle$ (nm)&0,27694&0,15384\\
Distanza minima atomi pesanti--G (nm)&0,51385&0,73865\\
Interazione diretta P--G (kJ/mol)&$-5,97$&$-0,85$\\
Legami H HFIP--carbonile&3,25&3,54\\
\bottomrule
\end{tabular}
\end{table}

La catena proveniente dallo stato aderente resta quasi verticale;
un'estremit\`a \`e pi\`u vicina al foglio. La seconda resta
quasi parallela. La solvatazione del carbonile non separa da sola
i due stati. Il risultato fornisce una variabile concreta per il
campionamento, anzich\'e attribuire il mancato distacco a una
generica scarsa solubilit\`a.

Se $F(z,q)$ comprende la misura statistica delle due coordinate,
il profilo lungo $z$ soddisfa, a meno di una costante,
\begin{equation}
 F(z)=-RT\log\!\int e^{-F(z,q)/(RT)}\,\mathrm dq+C.
\end{equation}
Una traiettoria che mantiene una sola popolazione in $q$ non
esegue questa integrazione, anche se le distribuzioni di $z$
sembrano collegabili. Abbiamo implementato $R_{g,z}$ nel motore
di simulazione e verificato i valori su 1.002 configurazioni:
errore massimo inferiore a $4\times10^{-10}$\,nm rispetto al
calcolo indipendente. Le forze del vincolo concordano con la
formula analitica entro $5\times10^{-7}$ e con le derivate
numeriche entro $3{,}7\times10^{-4}$\,kJ\,mol$^{-1}$\,nm$^{-1}$.
Due nuovi campioni di 100\,ps con vincoli simultanei su $z$ e
$R_{g,z}$ rimangono distinti. Il controllo numerico \`e riuscito;
il campionamento resta insufficiente e il profilo non viene accettato.

\subsection{Un generatore polarizzabile di parametri e il limite della verifica nel vuoto}

Abbiamo eseguito ByteFF-Pol \cite{ByteFFPol2026}, che predice
parametri molecolari con una rete e poi usa un potenziale esplicito
con polarizzazione. Sono stati generati sette insiemi di parametri
e verificati tredici complessi sulle stesse geometrie degli altri
confronti. Le energie sono state scomposte e ricomposte;
la forza netta sul frammento coincide con la derivata numerica
dell'energia entro $2,3\times10^{-6}$\,kJ\,mol$^{-1}$\,\AA$^{-1}$
nel controllo HFIP--acetato di metile. A grande distanza
l'interazione tende a zero.

L'interazione calcolata \`e $-34,57$\,kJ/mol per HFIP--acetato,
$-28,81$ per acido acetico--acetato e $-20,86$ per IPA--acetato.
Sono valori diversi dai riferimenti PBE0 della Tabella~\ref{tab:quantum}
e da MACE-POLAR-1-M, che d\`a rispettivamente $-46,54$, $-45,10$
e $-27,15$. Il corretto controllo delle forze non risolve
questa incertezza fisica. Abbiamo quindi eseguito dieci calcoli
autoconsistenti con $\omega$B97M-V/def2-TZVPD, lo stesso livello
usato per le interazioni di addestramento, includendo VV10 senza D3.
Per HFIP--acetato il valore corretto \`e $-46,165$\,kJ/mol;
il passaggio fra due griglie non locali cambia il risultato
di 0,071\,kJ/mol. Su questa geometria MACE-POLAR-1-M differisce
di 0,37\,kJ/mol, mentre ByteFF-Pol sottostima l'attrazione
di circa 11,6\,kJ/mol. Il controllo seleziona M per proseguire
lo studio di questa interazione locale e sospende l'uso predittivo
di ByteFF per HFIP prima di correggerlo o verificarlo su altre
geometrie. Non \`e una validazione universale di M, n\'e
un risultato sull'efficacia di HFIP come pulente.

Il modello non include silicio. La sua esecuzione nel vuoto
non dimostra densit\`a, permittivit\`a o desorbimento corretti.
L'esportazione verso OpenMM richiede inoltre le scale elettrostatiche
intramolecolari specifiche della pubblicazione; la versione installata
non \`e stata modificata globalmente e non viene usata come
se fosse gi\`a quel motore. Questa \`e una via da qualificare,
non una soluzione dedotta dal nome del programma.

\subsection{Esecuzione del modello polarizzabile nel liquido sul Mac}

Abbiamo reso eseguibile MACE-POLAR-1-M su una cella periodica
di 64 molecole HFIP, 768 atomi, sulla GPU del Mac disponibile.
I pesi del modello sono preservati. Il calcolo delle reti radiali,
dei prodotti e delle somme sui vicini procede a blocchi;
le derivate vengono ricalcolate dai rispettivi ingressi.
Le operazioni sono equivalenti per le prime derivate, con
parametri congelati. Non viene usato questo percorso per Hessiane
o addestramento.

Una contrazione onerosa pu\`o inoltre essere riordinata senza
costruire tutte le coppie di canali:
\begin{equation}
 \sum_{v,d}w_{uv}x_{bud}y_{bvd}
 =\sum_d x_{bud}\left(\sum_v w_{uv}y_{bvd}\right).
\end{equation}
La sottrazione dei riferimenti atomici additivi \`e una scelta
dello zero di energia che preserva forze e differenze a composizione
fissa. Una cella periodica diluita di 96 atomi controlla il percorso
contro il calcolo originale in doppia precisione: le modifiche
al riordino delle operazioni preservano le forze entro
$2\times10^{-15}$\,eV/\AA. Il confronto con la GPU in precisione singola
ha scarto relativo quadratico medio $1{,}5\times10^{-5}$.
Sulla cella completa il confronto processore--GPU, entrambi in
precisione singola, d\`a $1{,}1\times10^{-6}$; lo scarto energetico
\`e $1{,}3\times10^{-4}$\,eV. Il controllo completo in doppia
precisione sul processore supera il limite di memoria imposto
e non viene considerato concluso.

La GPU \`e limitata al 25\% della memoria assegnabile al processo;
il calcolo effettivo delle forze richiede circa 9,9 secondi
per configurazione. Le derivate numeriche di energia rispetto
a spostamento atomico, traslazione di una molecola e deformazione
della cella concordano con le derivate calcolate entro lo 0,34\%,
1,24\% e 0,84\% ai passi documentati. Due traiettorie a energia
totale costante, con identica partenza e durata di 8\,fs, hanno
escursioni energetiche massime di 0,03929 e 0,00979\,eV per
passi di 0,5 e 0,25\,fs. La riduzione \`e coerente con l'errore
quadratico del metodo di integrazione.

Queste verifiche rendono concretamente disponibile un modello
pi\`u avanzato per ulteriori controlli. La configurazione iniziale
proviene dal modello classico; otto femtosecondi non misurano
densit\`a o permittivit\`a di equilibrio. La qualificazione del
liquido e dell'interfaccia rimane un passaggio distinto.

\subsection{Pressione coerente con il potenziale: una correzione necessaria}

Nel percorso installato del modello polarizzabile la deformazione
virtuale della cella modificava i vettori locali, mentre posizioni,
cella reciproca e volume usati dalla parte elettrostatica restavano
non deformati. Questo non altera l'energia a deformazione nulla,
ma omette contributi alla sua derivata. Abbiamo quindi verificato
la pressione in doppia precisione su una configurazione di
16 molecole HFIP, confrontandola con
\begin{equation}
 \sigma_{ii}=\frac{1}{V}\frac{\partial U}{\partial\epsilon_{ii}},
 \qquad P_{\rm pot}=-\frac{\operatorname{tr}\sigma}{3}.
\end{equation}
Lo scarto isotropo \`e circa 183,56\,bar e persiste dimezzando
il passo di deformazione: non \`e un errore della precisione
singola. La prima traiettoria a cella variabile viene quindi
conservata come tentativo rifiutato, senza ricavarne propriet\`a
del liquido.

Abbiamo reso differenziabili anche le geometrie elettrostatiche
deformate e i fattori reciproci gaussiani. Le operazioni a
deformazione nulla mantengono esattamente l'energia e cambiano
le forze di meno di $1{,}5\times10^{-15}$\,eV/\AA\ nel confronto
in doppia precisione. Lo scarto residuo dalla derivata numerica
\`e 0,0106\,bar al passo $5\times10^{-5}$, con riduzione quadratica
al dimezzamento del passo. I pesi appresi sono preservati.

Ripetendo la stessa partenza e lo stesso integratore a cella
variabile, negli iniziali 50\,fs l'escursione massima dell'energia
conservata estesa passa da 0,04393 a 0,001056\,eV: una riduzione
di circa 42 volte. Questo verifica la correzione fisica del
calcolo. La nuova traiettoria di 100\,fs rimane transitoria:
la densit\`a scende dalla condizione iniziale 1,607 a circa
1,236\,g/cm$^3$. Non \`e una densit\`a di equilibrio accettata.

\subsection{Campionare le orientazioni senza attendere la dinamica lenta}

Il calcolo della sola energia del medesimo modello richiede
circa 0,77 secondi per una configurazione da 192 atomi;
il controllo contro il percorso con forze differisce al massimo
di $7{,}7\times10^{-6}$\,eV nei tre punti verificati.
Abbiamo quindi eseguito un campionamento a pressione imposta
che combina traslazioni, rotazioni molecolari, rotazioni del
legame OH e brevi proposte di dinamica hamiltoniana. Tutte le
coordinate interne restano aggiornabili; non si identifica
il liquido con molecole permanentemente rigide.

Per una proposta simmetrica in $\log V$, che scala solo i
centri molecolari conservando le coordinate relative,
l'accettazione \`e
\begin{equation}
 \min\left\{1,\exp\left[-\beta(\Delta U+P\Delta V)
 +(N_{\rm mol}+1)\log\frac{V'}{V}\right]\right\}.
\end{equation}
Il fattore $N_{\rm mol}$ viene dal determinante della
trasformazione dei centri; l'unit\`a aggiuntiva deriva dalla
proposta in volume logaritmico. Per 16 centri ideali, con
$\beta P=1$, il volume deve avere media e varianza 17.
Il controllo indipendente restituisce rispettivamente 16,993
e 17,008. Le proposte hamiltoniane usano quattro passi di
0,25\,fs e vengono accettate attraverso la variazione
dell'energia totale; il numero di mosse non misura un tempo
fisico di pulizia.

La prima serie effettiva di 1.000 mosse comprende
traslazioni, rotazioni, torsioni OH e 19 proposte hamiltoniane.
Le densit\`a campionate scendono fino a circa 1,268\,g/cm$^3$;
non si osservano rifiuti dovuti ai limiti di volume imposti.
L'equilibrio, le partenze indipendenti e l'effetto della
dimensione della cella restano da controllare. Il campionamento
continua preservando configurazione, contatori e stato del
generatore casuale. Non adattiamo il modello alla densit\`a
sperimentale per dichiarare riuscito il confronto.

\subsection{Un controllo chimico ulteriore: autoassociazione dei solventi}

Una simile attrazione locale per il carbonile non implica
una simile competizione nel solvente. Abbiamo aggiunto i
dimeri HFIP--HFIP e acido acetico--acido acetico, mantenendo
lo stesso livello $\omega$B97M-V/def2-TZVPD e la correzione
di base. Sulla griglia non locale iniziale otteniamo interazioni
congelate di $-29{,}744$ e $-88{,}060$\,kJ/mol. L'affinamento
HFIP porta a $-29{,}932$, con variazione di griglia 0,187\,kJ/mol;
MACE-POLAR-M restituisce $-29{,}984$, scarto circa 0,052\,kJ/mol.
Per il dimero dell'acido la griglia affinata d\`a $-88{,}074$,
variazione 0,014\,kJ/mol; il modello d\`a $-88{,}924$. Il dimero dell'acido
ha due legami a idrogeno; il confronto non li tratta come un
singolo legame equivalente. Il dimero non qualifica il liquido intero.

Per evitare di confondere interazione e deformazione,
abbiamo anche costruito, nel medesimo modello, la reazione
elettronica nel vuoto
\begin{equation}
 S_2+\mathrm{acetato}\longrightarrow S\!:\!\mathrm{acetato}+S,
\end{equation}
usando gli stessi monomeri raffinati come riferimento comune.
L'energia risultante \`e circa $-8,57$\,kJ/mol per HFIP e
$+32,75$\,kJ/mol per l'acido acetico. Questo individua una
competizione chimica diversa, assente dal criterio basato sul
solo legame solvente--estere. \`E una reazione fra piccoli
aggregati: entropia, molteplici partner nel liquido e interfaccia
sono ancora assenti. Inoltre la solvatazione del carbonile
pu\`o gi\`a essere presente sia nella catena aderente sia in
quella libera. Il risultato sostiene un'ipotesi da approfondire,
non seleziona da solo un bagno che pulisca il campione.

\subsection{Vie indipendenti e lettura critica delle nuove proposte}

Abbiamo tentato anche la teoria statistica del liquido DRISM/KH,
imponendo la permittivit\`a sperimentale. I modelli organici
non raggiungono il criterio di convergenza, neppure nei tentativi
documentati di variazione del solutore e prosecuzione in densit\`a
e carica. L'acqua di controllo distribuita con il programma
converge in 27 iterazioni. Non usiamo suscettivit\`a o energie
di solvatazione organiche non convergenti.

Il trattamento con nitrito del 2025 \cite{LeeNitrite2025} avviene
prima del deposito, lascia particelle e riduce il rapporto XPS
C--O--C da 20\% a 17\%; non prova il salvataggio richiesto.
Un trattamento elettrico senza bolle \cite{Park2018} \`e invece
applicato dopo il trasferimento. Richiede per\`o contatto elettrico
al grafene e un elettrolita: resta fuori dalla pulizia con soli bagni
e non viene presentato come immediatamente compatibile con
cristalli isolati da prelevare. Un miglioramento elettrico deve
comunque essere distinto dalla scomparsa del polimero.

\textbf{La ricerca computazionale \`e attiva. Restano aperti il
campionamento interfaciale, la qualificazione dei solventi,
PPC e lunghezze maggiori, il substrato reale e la verifica sul campione.}
I risultati disponibili non autorizzano una classifica dei pulenti
o una dichiarazione di soluzione.

\section{Decisione scientifica e prossimo risultato da cercare}

La nuova proposta consiste nel \emph{modificare l'ordine e la
tempistica di una sequenza gi\`a parzialmente attiva}, controllando
quando il residuo diventa estraibile e quando torna aderente.
Il risultato teorico \`e che questo effetto \`e possibile, ha
condizioni di segno e pu\`o essere distinto da una semplice
somma di tempi in solventi. Le simulazioni escludono una
promessa generale di efficacia dei cicli.

Se I migliora rispetto a R senza danno, il primo progresso
operativo \`e una sequenza DMF--THF senza IPA intermedio.
Se l'IPA \`e utile e il confronto alternato supera quello
a blocchi, il progresso \`e un'estrazione a cicli da ottimizzare
con un bilancio di massa. Se la massa trascinabile non pu\`o
spiegare il residuo, l'asciugatura perde priorit\`a. Se il
residuo ha un'identit\`a diversa, il bersaglio chimico va
corretto prima di scegliere altri reagenti.

Nessuna di queste conclusioni \`e stata ancora osservata sul
campione di Armando. Il contributo qui ottenuto \`e una teoria
con previsioni controllabili, calcoli riproducibili e due
interventi nuovi rispetto ai tentativi riferiti. Non viene
rivendicata priorit\`a scientifica della strategia o risoluzione
sperimentale del problema.

Il nuovo risultato da cercare \`e una differenza di energia libera
e di barriera che resista ai controlli di modello e campionamento,
insieme a una finestra di selettivit\`a sul substrato reale. La scelta
dei bagni deve poi superare le prove di massa, topografia, chimica
e resa del prelievo gi\`a specificate. La campagna molecolare non
ha ancora raggiunto questo risultato.

\appendix
\section{Verifica dei calcoli e riproduzione}

Il programma \texttt{model\_cleaning.py} richiede Python e NumPy.
Il comando
\begin{center}
\texttt{python model\_cleaning.py --out-dir risultati}
\end{center}
genera \texttt{risultati.json} e le coordinate dei grafici
in \texttt{figure-incorporate.tex}. I grafici di questo
documento contengono direttamente tali coordinate e non
dipendono da immagini esterne.

Sono stati controllati 540 casi del modello di trasporto
(6 valori di $\beta$, 6 di $r$, 3 di $\varepsilon$ e 5 tempi).
L'errore massimo di conservazione della matrice di transizione
\`e $7{,}8\times10^{-14}$; lo scarto dalla soluzione analitica
del bagno senza ricambio \`e $2{,}0\times10^{-14}$.
Il minimo elemento delle matrici \`e zero. Sono inoltre
verificati il limite di estrazione ideale, la catena
irreversibile $L\to M\to R$, l'invarianza quando i due
generatori coincidono e il limite dei ricambi a risorse fisse.
Un'integrazione indipendente a passi piccoli conferma
l'evoluzione anche con popolazione inizialmente schermata,
con scarto inferiore a $10^{-12}$ nei casi confrontati.

Questi controlli attestano la soluzione delle equazioni
implementate. Non attestano che le equazioni descrivano
correttamente il residuo reale. Le costanti $a,b,e$, $k_d$,
$k_a$, $k_m$ e la frazione inaccessibile restano da identificare
con esperimenti. Non \`e stata eseguita una simulazione
atomistica del sistema PMMA950K/PPC/grafene/\SiO\ nei
solventi effettivi.

\section{Derivazioni brevi dei limiti usati}

Per la formula~\eqref{eq:sign}, durante P la massa mobile
varia di
\[
 M(t)-M_0=\frac{aL_0-bM_0}{a+b}
                 (1-e^{-(a+b)t}).
\]
Durante G viene esportata la frazione
$k(1-e^{-(k+c)u})/(k+c)$ della massa mobile iniziale.
Il prodotto d\`a la differenza fra i due ordini.

Per~\eqref{eq:resource}, in ogni bagno la quantit\`a
$\beta$ diventa $NB$ e il tempo $\tau_{\rm tot}/N$.
Si applica~\eqref{eq:static} e si moltiplicano i fattori
perch\'e ogni nuovo bagno riparte pulito. Il logaritmo del
fattore \`e asintotico a
$-(1-e^{-B\tau_{\rm tot}})/(NB)$, da cui
\eqref{eq:resource-limit}. Lasciare un bagno carico non
equivale a questo azzeramento della concentrazione.

Per~\eqref{eq:counterexample}, in ciascuna fase D la massa
$L$ si moltiplica per $e^{-w/N}$; nella fase T accade
lo stesso a $M$. Sommando le quote mobilizzate nei vari
stadi si ottiene la formula. La funzione
$N(1-e^{-w/N})$ cresce strettamente: ponendo $x=w/N>0$,
la sua derivata \`e $1-(1+x)e^{-x}>0$.
Quindi in questo caso la rimozione peggiora con pi\`u cicli.


\begin{thebibliography}{99}
\bibitem{Allresist}
Allresist GmbH, \emph{AR-P 630--670: PMMA-Resists}, scheda tecnica,
serie AR-P 672, prodotto 672.045.
\href{https://www.allresist.de/wp-content/uploads/2020/03/AR-P630-670_Deutsch_Allresist_Produktinformation.pdf}{Scheda del produttore}.

\bibitem{Tyagi2022}
A. Tyagi, V. Mi\v{s}eikis, L. Martini et al.,
\emph{Ultra-clean high-mobility graphene on technologically relevant substrates},
Nanoscale \textbf{14} (2022), 2167--2176.
\href{https://doi.org/10.1039/D1NR05904A}{doi:10.1039/D1NR05904A}.
\href{https://www.rsc.org/suppdata/d1/nr/d1nr05904a/d1nr05904a1.pdf}{Supplemento consultato}.

\bibitem{Xiao2019}
Z. Xiao, Q. Wan e C. Durkan,
\emph{Cleaning Transferred Graphene for Optimization of Device Performance},
Advanced Materials Interfaces \textbf{6} (2019), 1801794.
\href{https://doi.org/10.1002/admi.201801794}{doi:10.1002/admi.201801794}.
\href{https://www.repository.cam.ac.uk/handle/1810/291256}{Manoscritto degli autori consultato}.

\bibitem{Manjkow1987}
J. Manjkow, J. S. Papanu, D. S. Soong, D. W. Hess e A. T. Bell,
\emph{An in situ study of dissolution and swelling behavior of
poly-(methyl methacrylate) thin films in solvent/nonsolvent binary mixtures},
Journal of Applied Physics \textbf{62} (1987), 682--688.
\href{https://doi.org/10.1063/1.339742}{doi:10.1063/1.339742}.
Consultato il riassunto primario, non il testo integrale.

\bibitem{Rooks2002}
M. J. Rooks et al., \emph{Low stress development of poly(methylmethacrylate)
for high aspect ratio structures}, Journal of Vacuum Science \& Technology B
\textbf{20} (2002), 2937--2941.
\href{https://doi.org/10.1116/1.1524971}{doi:10.1116/1.1524971}.
\href{https://nano.yale.edu/sites/default/files/files/pmma_develop_ipa_water.pdf}{Testo consultato}.

\bibitem{Huston2020}
K. J. Huston, C. E. Rice e R. G. Larson,
\emph{Forward Flux Sampling of Polymer Desorption Paths from a Solid Surface
into Dilute Solution}, Polymers \textbf{12} (2020), 2275.
\href{https://doi.org/10.3390/polym12102275}{doi:10.3390/polym12102275}.

\bibitem{Ayodele2022}
O. O. Ayodele, S. Pourianejad, A. Trofe, A. Prokofjevs e T. Ignatova,
\emph{Application of Soxhlet Extractor for Ultra-clean Graphene Transfer},
ACS Omega \textbf{7} (2022), 7297--7303.
\href{https://doi.org/10.1021/acsomega.1c07113}{doi:10.1021/acsomega.1c07113}.

\bibitem{Tang2025}
Z. Tang et al., \emph{A Closed-Loop Solvent Recycling Device for Polymer
Removal in Graphene Transfer Process}, Separations \textbf{12} (2025), 295.
\href{https://doi.org/10.3390/separations12110295}{doi:10.3390/separations12110295}.

\bibitem{Duan2022}
T. Duan, H. Li, R. Papadakis e K. Leifer,
\emph{Towards ballistic transport CVD graphene by controlled removal of
polymer residues}, Nanotechnology \textbf{33} (2022), 495704.
\href{https://doi.org/10.1088/1361-6528/ac8d9b}{doi:10.1088/1361-6528/ac8d9b}.

\bibitem{Merino2024}
J. P. Merino, S. Brosel-Oliu, G. Rius et al.,
\emph{Ethanol Solvation of Polymer Residues in Graphene Solution-Gated
Field Effect Transistors}, ACS Sustainable Chemistry \& Engineering
\textbf{12} (2024), 9133--9143.
\href{https://doi.org/10.1021/acssuschemeng.4c01538}{doi:10.1021/acssuschemeng.4c01538}.

\bibitem{Suk2013}
J. W. Suk et al., \emph{Enhancement of the Electrical Properties of Graphene
Grown by Chemical Vapor Deposition via Controlling the Effects of Polymer Residue},
Nano Letters \textbf{13} (2013), 1462--1467.
\href{https://doi.org/10.1021/nl304420b}{doi:10.1021/nl304420b}.

\bibitem{Schoenemann2018}
E. Sch\"onemann, A. Laschewsky e A. Rosenhahn,
\emph{Exploring the Long-Term Hydrolytic Behavior of Zwitterionic
Polymethacrylates and Polymethacrylamides}, Polymers \textbf{10} (2018), 639.
\href{https://doi.org/10.3390/polym10060639}{doi:10.3390/polym10060639}.

\bibitem{Jung2006}
J. H. Jung, M. Ree e H. Kim,
\emph{Acid- and base-catalyzed hydrolyses of aliphatic polycarbonates and
polyesters}, Catalysis Today \textbf{115} (2006), 283--287.
\href{https://doi.org/10.1016/j.cattod.2006.02.060}{doi:10.1016/j.cattod.2006.02.060}.
Consultati riassunto e passaggio metodologico accessibile.

\bibitem{Wang2017}
X. Wang et al., \emph{Direct Observation of Poly(Methyl Methacrylate)
Removal from a Graphene Surface}, Chemistry of Materials
\textbf{29} (2017), 2033--2039.
\href{https://doi.org/10.1021/acs.chemmater.6b03875}{doi:10.1021/acs.chemmater.6b03875}.

\bibitem{Schwartz2019}
J. J. Schwartz et al., \emph{Chemical Identification of Interlayer
Contaminants within van der Waals Heterostructures},
ACS Applied Materials \& Interfaces \textbf{11} (2019), 25578--25585.
\href{https://doi.org/10.1021/acsami.9b06594}{doi:10.1021/acsami.9b06594}.

\bibitem{Arai1996}
T. Arai, N. Sawatari, T. Yoshizaki, Y. Einaga e H. Yamakawa,
\emph{Excluded-Volume Effects on the Hydrodynamic Radius of Atactic and
Isotactic Oligo- and Poly(methyl methacrylate)s in Dilute Solution},
Macromolecules \textbf{29} (1996), 2309--2314.
\href{https://doi.org/10.1021/ma951274n}{doi:10.1021/ma951274n}.
Parametro impiegato dalla Tabella 3.

\bibitem{WangHFIP2020}
Y. Wang, T. Duo, X. Xu et al., \emph{Eco-Friendly High-Performance
Poly(methyl methacrylate) Film Reinforced with Methylcellulose},
ACS Omega \textbf{5} (2020), 24256--24261.
\href{https://doi.org/10.1021/acsomega.0c02249}{doi:10.1021/acsomega.0c02249}.

\bibitem{Marchelli2022}
G. Marchelli, J. Ingenmey, O. Holl\'oczki, A. Chaumont e B. Kirchner,
\emph{Hydrogen Bonding and Vaporization Thermodynamics in
Hexafluoroisopropanol-Acetone and -Methanol Mixtures.
A Joined Cluster Analysis and Molecular Dynamic Study},
ChemPhysChem \textbf{23} (2022), e202100620.
\href{https://doi.org/10.1002/cphc.202100620}{doi:10.1002/cphc.202100620}.
Parametri dal supplemento consultato integralmente.

\bibitem{Casoria2024}
M. Casoria, M. Macchiagodena, P. Rovero et al.,
\emph{Upgrading of the general AMBER force field 2 for fluorinated
alcohol biosolvents: A validation for water solutions and melittin solvation},
Journal of Peptide Science \textbf{30} (2024), e3543.
\href{https://doi.org/10.1002/psc.3543}{doi:10.1002/psc.3543}.

\bibitem{Behbahani2021}
A. Foroozani Behbahani e V. Harmandaris,
\emph{Gradient of Segmental Dynamics in Stereoregular Poly(methyl
methacrylate) Melts Confined between Pristine or Oxidized Graphene Sheets},
Polymers \textbf{13} (2021), 830.
\href{https://doi.org/10.3390/polym13050830}{doi:10.3390/polym13050830}.
Consultati Tabelle S1--S2 e correzione
\href{https://doi.org/10.3390/polym14010106}{doi:10.3390/polym14010106}.

\bibitem{MACE2025}
D. P. Kov\'acs, J. H. Moore, N. J. Browning et al.,
\emph{MACE-OFF: Transferable Short Range Machine Learning Force Fields
for Organic Molecules}, arXiv:2312.15211, versione 5 (2025).
\href{https://arxiv.org/abs/2312.15211v5}{Manoscritto};
\href{https://github.com/ACEsuit/mace-off}{Modello ufficiale}.

\bibitem{MACEPolar2026}
I. Batatia, W. J. Baldwin, D. Kuryla et al.,
\emph{MACE-POLAR-1: A Polarisable Electrostatic Foundation Model
for Molecular Chemistry}, arXiv:2602.19411 (2026).
\href{https://arxiv.org/abs/2602.19411}{Manoscritto};
\href{https://github.com/ACEsuit/mace-foundations/releases/tag/mace_polar_1}{Modelli ufficiali}.

\bibitem{ByteFFPol2026}
T. Zheng, X. Xu, Z. Wang et al.,
\emph{Bridging quantum mechanics to liquid properties via a universal
organic force field}, Nature Communications \textbf{17} (2026), 8276.
\href{https://doi.org/10.1038/s41467-026-73566-3}{DOI: 10.1038/s41467-026-73566-3}.

\bibitem{LeeNitrite2025}
K. Lee, J. Kil, J. Park, S. Yang e B. Park,
\emph{Rapid and Efficient Polymer/Contaminant Removal from Single-Layer
Graphene via Aqueous Sodium Nitrite Rinsing for Enhanced Electronic Applications},
Polymers \textbf{17} (2025), 689.
\href{https://doi.org/10.3390/polym17050689}{DOI: 10.3390/polym17050689}.

\bibitem{Park2018}
B. Park, J. N. Huh, W. S. Lee e I.-G. Bae,
\emph{Simple and rapid cleaning of graphenes with a `bubble-free'
electrochemical treatment}, Journal of Materials Chemistry C
\textbf{6} (2018), 2234--2244.
\href{https://doi.org/10.1039/C7TC05695H}{DOI: 10.1039/C7TC05695H}.

\bibitem{Hunger2010}
J. Hunger, R. Buchner, M. E. Kandil, E. F. May, K. N. Marsh e G. Hefter,
\emph{Relative Permittivity of Dimethylsulfoxide and N,N-Dimethylformamide
at Temperatures from (278 to 328) K and Pressures from (0.1 to 5) MPa},
Journal of Chemical \& Engineering Data \textbf{55} (2010), 2055--2065.
\href{https://doi.org/10.1021/je9010773}{DOI: 10.1021/je9010773}.

\end{thebibliography}
\end{document}
